% Performance model and throughput results.
\section{Performance analysis}
\label{sec:performance}
This analysis asks how much communication and execution work each lock
requires, and how that work constrains throughput and caller-return latency.
The service policy, executed operation mix, and executor sequence are inputs;
no equal-share or bounded-lag guarantee is assumed. Results here do not use
the fairness theorems. Costs are elapsed time, resource occupancy, or source
work counts, not a scheduler's service-credit metric.

\subsection{Communication and execution model}
\label{sec:model}
\subsubsection{Communication events}
Classical LogP~\cite{logp} specifies small-message latency bound $L$, processor
send/receive occupancy $o$, endpoint send/receive gap $g$, and $P$ processor
modules. Its capacity is at most $\lceil L/g\rceil$ messages in flight from
or to an endpoint. The runtime convention charges $L$ per message, giving
$L+2o$ for an isolated one-way delivery. Independent deliveries can overlap:
latency belongs to dependency paths, while occupancy and gaps constrain rate.
An upper-bound $L$ is not a physical lower-bound latency.

For coherent shared memory we use a \emph{LogP-inspired adaptation}, not an
identification of a load with a message. Repeated private-cache polling can
produce no new fabric traffic. Ownership transfer can require requests,
invalidations, acknowledgments, and forwarding. Following the motivation of
cache-coherent communication models~\cite{coherence}, transaction costs are
indexed by type and topology. $P$ counts hardware execution endpoints and
$N$ counts clients; directory/interconnect bottlenecks may be additional
resources. A large cache line is not automatically a LogGP long message:
LogGP's byte-gap parameter $G$~\cite{loggp} is appropriate only when a calibrated
bulk-transfer regime exists.


Construct a finite event DAG containing delegate execution, publication,
discovery, queue manipulation, transfer, completion, and executor changes.
Delegate nodes touching the same protected state are serialized. Independent
publication and notification may overlap this chain. Label CPU work with its
endpoint, messages with endpoint/topology, and contended atomic ownership
with the resource it occupies. Failed elections, empty scans, and parked
threads have events but contribute no successful operations.

\begin{proposition}[Resource and dependency bound]
\label{prop:resource}
For a closed finite DAG completing $B$ operations, let $\mathrm{CP}$ be its
longest weighted path, $W_p$ its total non-overlappable CPU occupancy at
endpoint $p$, $v_p^s,v_p^r$ its message-send/receive counts, and $r_\ell a_\ell$
the occupancy of a serialized hot-line resource. With endpoint gap $g_p$,
\begin{align}
T &\ge \max\{\mathrm{CP},\max_p W_p,\max_p G_p,
                    \max_\ell r_\ell a_\ell\},\label{eq:resource}\\
G_p&=\max\{(v_p^s-1)_+,(v_p^r-1)_+\}g_p .\nonumber
\end{align}
The costs must be exact within the abstract machine or guaranteed minima.
Capacity and any additional fabric constraints remain part of the DAG/model.
\end{proposition}
\begin{proof}
Every dependency path must complete; each non-overlappable resource must
supply its total occupancy. A sequence of $v$ sends or receives has $v-1$
inter-event gaps. Each requirement separately lower-bounds elapsed time;
their maximum is therefore a lower bound. Taking a sum would unjustifiably
forbid overlap. Include startup/drain events rather than counting transfers
that lie outside the chosen observation interval.
\end{proof}
Thus $B/T$ has a model upper bound, not necessarily an attainable rate.
Substituting measured \emph{means} for guaranteed minima yields a performance
estimate, not a hardware theorem. Moreover, $O(1)$ successful queue operations
do not imply $O(1)$ elapsed time when a predecessor can stop publishing.

\subsubsection{Trace accounting and a predictive approximation}
Let a trace have $E$ passes, $B>0$ completions, and $F$ executor changes.
Define $b=B/E$ and $p=F/E$; zero-completion passes count in $E$.
Partition the \emph{serialized execution lane}, without double counting,
into warm service, per-operation scheduling/discovery and exposed request
communication, pass administration, exposed protected-state/scheduler-state
reacquisition, and uncovered idle time. Mean per-completion time is then
represented by
\begin{equation}
 t_D \simeq \bar C_D+d_D+
 \frac{A_D+p_D(M_D+K_D)}{b_D}+I_D .\label{eq:delegation}
\end{equation}
$A$ is mean administration per pass, $M$ and $K$ are mean exposed protected-data
and scheduler-state reacquisition \emph{conditional on an executor change},
and $I$ is uncovered idle time per completion. When $F=0$, its migration term
is zero; finite startup cost is separate. $\bar C$ and $d$ are weighted by
operations, not passes. Ratios of totals give this accounting; averaging
per-pass ratios or assuming independence of batch size and cost does not.

As a trace partition, the accounting is definitional. As a prediction,
Equation~\eqref{eq:delegation} replaces trace costs by independently calibrated
values and is only an approximation. It does not replace the DAG constraints.
A request miss inside measured service must not be added again to $d$;
a cold protected-data miss must not be counted both in $C$ and $M$.
Speculative/prefetched communication contributes only its \emph{exposed}
residual to this serialized decomposition. Network demand still appears in
Equation~\eqref{eq:resource} even if its latency was hidden.

Protected-data cost $M(W,\tau,\rho)$ depends on touched reusable lines $W$,
topology $\tau$, and residency $\rho$, not merely allocated object size.
Independent transfers might fit $d_\tau+(W-1)\gamma_\tau$; a dependent
pointer chain might fit $W d_\tau$. These are alternative calibrated regimes,
not universal coherence formulas. Input lines, queue metadata and protected
state are separate working sets. A large heap can cease to be cache resident.
An unchanged executor can also suffer capacity misses or interference;
those costs are not excluded by the absence of executor-change migration.

A comparable handoff estimate is
\begin{equation}
 t_H\simeq \bar C_H+h_H+M_H,\label{eq:handoff}
\end{equation}
where $h_H$ includes lock protocol and policy work, and $M_H$ is actual exposed
movement per operation, possibly zero for local reuse. Admission cooldowns
are not automatically $I_D$: another request may use the server meanwhile.

\subsection{Migration amortization and throughput crossover}
\label{sec:amortization}
\begin{proposition}[Intra-epoch ownership invariance]
\label{prop:ownership}
Fix an epoch's executor, its operations' accessed-line trace up to a
permutation, and exclusive protected-state access by that executor during the
epoch. Permuting request order cannot introduce a transfer of a protected
line to a \emph{different executing core within that epoch}.
\end{proposition}
\begin{proof}
Every protected access in the permuted epoch is still issued by the same
core. There is no second executor to which execution ownership can transfer.
This says nothing about first-touch acquisition from an earlier owner,
eviction/refetch from memory, requester inputs, or off-path metadata traffic.
\end{proof}
This is narrower than ``the data always stays in L1.'' If ordering changes
which lines are accessed, the initial owners of those lines, the operation
mix, batch boundaries, or executor selection, even the migration comparison
needs to be recomputed.

\begin{corollary}[Resident repeated-footprint case]
For a fixed $W$-line exclusively modified footprint touched in every epoch,
no external invalidations/evictions, and $F$ executor changes, at most $WF$
protected-line ownership changes are required after warmup; under the model
where every changed executor must acquire every line, the count is exactly
$WF$. Per completion it is $Wp/b$.
\end{corollary}
\begin{proof}
An unchanged executor requires no reacquisition. Each changed executor
acquires each of the $W$ lines once, then retains it for that epoch.
Sum over changed epochs and divide by $B$, using $F/B=p/b$.
\end{proof}
Handoff locks correspond to one operation per epoch, but their switch
probability and topology must be measured. Owner-to-owner movement already
exists in MCS; it is not a cost introduced only by CFL's added policy.

\begin{proposition}[Conditional batching crossover]
\label{prop:crossover}
Assume the same service mix and warm cost $C$, the same migration cost $M$,
$p_D=1$, no uncovered idle time, and costs independent of batch size in the
considered regime. Within Equations~\eqref{eq:delegation}--\eqref{eq:handoff},
delegation has smaller time per operation exactly when
\begin{equation}
 b(M+h-d)>M+A+K .\label{eq:crossover}
\end{equation}
If $M+h-d>0$, the smallest positive integer completion batch satisfying this
inequality is
\begin{equation}
 b_{\min}=\left\lfloor\frac{M+A+K}{M+h-d}\right\rfloor+1.
\end{equation}
If the denominator is nonpositive and all costs are nonnegative, no positive
batch meets the strict inequality.
\end{proposition}
\begin{proof}
Cancel $C$ in $C+d+(M+A+K)/b<C+h+M$ and multiply by $b>0$.
The integer result follows from the strict inequality.
\end{proof}
The saved term is $(1-1/b)M$; the unpaid terms are per-request scheduling and
communication plus $(A+K)/b$. Increased migration cost lowers the relative
importance of scheduler work, but oversized inputs, a congested publication
line, or migrating PQ metadata can remove the advantage. This is a
regime-dependent prediction, not a ranking of all locks on all machines.

\subsubsection{Implementation deltas and controlled operation mix}
To explain the performance difference between two variants, compare:
\begin{align}
\Delta_D&\simeq\Delta \bar C+\Delta d+
 \Delta\!\left(\frac{A+p(M+K)}b\right)+\Delta I,\\
\Delta_H&\simeq\Delta\bar C+\Delta h+\Delta M.
\end{align}
FC-PQ versus FC changes discovery and batching as well as policy; CC-Ban
versus CC also changes the source budget (16 versus 64). CFL versus MCS
changes queue policy and locality, not the existence of handoff itself.
None is automatically a single-variable ablation. The evidence appendix
identifies these structural confounders before interpreting a slowdown.

For fixed per-class service $c_i>0$ and given useful-service fractions $f_i$
on a fully utilized overhead-free server,
\begin{equation}
 X_i=f_i/c_i,\qquad X=\sum_i f_i/c_i.\label{eq:mix}
\end{equation}
For costs 1 and 9, the supplied mix $(f_1,f_2)=(1/10,9/10)$ gives
$X=1/5$, whereas $(1/2,1/2)$ gives $X=5/9$ operations per normalized time
unit. No overhead changed. These fractions are inputs to the cost model,
not conclusions about which shares an implemented scheduler guarantees.
Report useful work, class completion rates, service shares, and total
operations together. A change in $\bar C$ due to a different class mix is not
necessarily synchronization inefficiency.

% Performance latency and implementation costs.
\subsection{Batching versus caller return latency}
\label{sec:return}
A waiter can return after its own notification, while an executor caller may
finish the rest of the batch first. If its operation is first in a batch of
exactly $b$ completed operations, each later operation with serialized elapsed
service at least $\tau_{\min}$ adds at least $(b-1)\tau_{\min}$
post-completion delay. This is an elapsed-time lower bound, excluding extra
administration. It follows from the synchronous return rule regardless of
the policy that selected the operation order.

Suppose a restricted implementation/regime has a certified per-subsequent-
operation serialized upper bound $\tau_{\max}$ and remaining fixed tail work
$e_{\rm tail}$. Then
\begin{equation}
 R_{\rm tail}\le e_{\rm tail}+(b-1)\tau_{\max}.\label{eq:tail}
\end{equation}
This assumes exactly $b$ completed operations and that \emph{all} other
post-completion work is in the envelope. It is not a bound for FC-PQ obtained
by replacing $b$ with 64: its announcement work, no-service pops, stalled
publishers, and retries also need bounds.

\begin{proposition}[Certified batching feasibility]
\label{prop:frontier}
Assume $\tau_{\max}>0$. Under
Proposition~\ref{prop:crossover}'s performance approximation and the
return envelope~\eqref{eq:tail}, a tail budget $D\ge e_{\rm tail}$ admits a
positive integer batch certified by both models iff
\begin{equation}
 b_{\min}\le b_{\max}:=
 1+\left\lfloor\frac{D-e_{\rm tail}}{\tau_{\max}}\right\rfloor,
 \qquad \tau_{\max}>0,
\end{equation}
with positive crossover denominator. If either condition fails, no batch
has both certificates under the assumed envelopes.
\end{proposition}
\begin{proof}
The strict throughput inequality requires $b\ge b_{\min}$.
Rearranging~\eqref{eq:tail} to meet $D$ gives $b\le b_{\max}$.
The integer interval is nonempty exactly under the stated condition.
\end{proof}
For an implementation with supported batches $\mathcal B$, further intersect
this integer interval with $\mathcal B$ and with the regime in which the
cost/envelope assumptions hold. A zero $\tau_{\max}$ is outside the stated
proposition; it would impose no batch-dependent tail limit when
$D\ge e_{\rm tail}$.
Absence of a certificate is not proof that real hardware cannot meet both
objectives: an upper envelope may be conservative. As a synthetic example,
let $(M,A,K,h,d)=(20,4,4,2,8)$, giving $b_{\min}=3$.
For $(e_{\rm tail},\tau_{\max},D)=(3,5,13)$, $b_{\max}=3$;
only $b=3$ is certified. Reducing $D$ to 8 gives $b_{\max}=2$ and no
certified batch. All numbers are normalized model costs, not measurements.
This locates the point where merely increasing a batch budget cannot satisfy
both modeled aims; a different return/executor mechanism or smaller
coordination cost is needed to improve the certified region.

For a stationary closed workload, the operation-weighted response-time law
$N=X(\mathbb E[R]+\mathbb E[Z])$, with $Z$ the outside-lock think time,
is an aggregate consistency check. It does not predict percentiles or an
individual request's maximum waiting time.

\subsection{Performance of the implemented lock families}
\label{sec:instances}
Table~\ref{tab:performance} summarizes the dominant \emph{source work}, not exact
fabric traffic. The appendix provides complete target coverage and source
ranges. Let $n$ be retained queue/list size, $a$ arrivals drained, $k$ pop
attempts, and $b$ actual completions; use $n_{\max}$ over a pass for complexity
bounds. Additional failed elections, allocation, stalls, and cleanup are
not assumed constant.

\begin{table*}[t]
\centering\small
\caption{Performance comparison only. Complexity describes source work,
not latency or one-message-per-instruction counts. Elapsed-time predictions
add machine, residency, and progress assumptions; no fairness theorem is used.}
\label{tab:performance}
\begin{tabular}{@{}P{0.14\textwidth}P{0.28\textwidth}P{0.25\textwidth}P{0.24\textwidth}@{}}
\toprule
Family & Discovery/scheduler demand & Execution and budget & Performance consequence \\
\midrule
FC & $O(n)$ visits/pass; $O(n/b)$ visits/completion for $b>0$ & Full list scan, no count cap; executor may change each pass & Retained idle nodes increase scan work; occupancy determines amortization. \\
FC-Ban & FC scan plus eligibility tests, skips and retained-registration accounting & Skips reduce $b$; uncovered idle differs from client cooldown & Charge scan visits and exposed idle, not every caller's entire cooldown. \\
FC-PQ heap/tree & $O((a+k)\log(n_{\max}+1))$ ordinary/amortized queue work plus buffers/timers & $a\le64$, $b\le k\le64$ per pass; queue may exceed 64 clients & Activation/pop ratios and migrating PQ metadata determine scheduling cost. \\
FC-SL & Ordered skip-set operations, allocation and atomic metadata costs & Count budget on pops, not completions; direct ordered publication & Shared ordered metadata adds traffic; current source has the shared correctness blocker. \\
CC / DSM & Linked-queue publication and $O(b)$ service traversal & CC cap 64, DSM source permits 65; both may stop earlier & Batch occupancy amortizes executor changes; absent links can stall progress. \\
CC-Ban & CC protocol plus pre-enqueue wait and charge accounting & Cap 16, unlike CC's 64 & Comparative cost includes budget difference and admission effects. \\
MCS / CLH & Constant successful queue/link operations; local/per-predecessor polling & Caller executes one operation per ownership epoch & Local polling does not remove protected-data migration or delayed linking. \\
Spin / Ticket & Central acquisition/turn cache lines; retry or sharer traffic & Caller execution; no delegation amortization & Contended line demand and placement matter beyond successful atomic counts. \\
Shfl / CFL & Variable queue inspection/relinking; CFL adds usage counters & Locality/policy can alter executor topology and run length & Model policy work and changed migration separately; neither is a fixed universal penalty. \\
Mutex / U-SCL & OS waiting or slice admission plus ownership protocol & Caller execution; kernel/preemption terms may dominate & Runtime implementation and scheduling assumptions are required for latency bounds. \\
\bottomrule
\end{tabular}
\end{table*}

\paragraph{FC versus FC-PQ.}
Under a full active FC scan, $b\approx n$ makes per-operation scan work
approximately constant. Under sparse activity, $n/b$ grows. In FC-PQ,
$k/b$ quantifies wasted completed-node pops, and $a/b$ activation intensity;
neither is fixed by $H$. Its scheduler ratio is approximately
$(a+k)\log(n_{\max}+1)/b$, not unconditionally $O(N\log N)$ per pass.
A heap's growth can copy $O(n)$ storage on a particular insertion, so an
amortized cost is not a per-pass worst-case latency guarantee. BinaryHeap
and BTreeSet share the ordering policy, not identical locality/allocation
costs; the shared PQ migrates when the combiner changes.

\paragraph{Queue and handoff families.}
CC/DSM amortize linked-queue service over an executor epoch. MCS/CLH
reduce centralized polling traffic relative to globally shared-turn polling,
but still have enqueue hotspots and owner-to-owner protected-state movement.
ShflLock and CFL change the topology and order of owners; the model therefore
changes both $M_H$ and policy work. Rust/C wrappers must be treated as
different implementations:
the C FC cleanup/polling rules and C CC batch limit are recorded in the appendix.


\paragraph{Admission and waiting strategies.}
FC-Ban scans and skips published requests; CC-Ban delays the requester before
enqueue. These paths change scan work, actual occupancy, and exposed idle
in different ways. U-SCL's slice admission needs its C protocol, not the
Rust wrapper alone. Mutex and blocking-Parker paths require wakeup and
scheduler terms; they are not accurately ranked by adding a fixed LogP
message latency to a spin-lock cost. A dedicated RCL executor, when enabled,
changes executor migration but consumes a server resource and scan capacity.
The checked-in dispatch limitations matter before claiming such a comparison.

\subsection{Performance validation and limits}
\label{sec:performance-validation}
A quantitative use of the model requires independently calibrated
transaction/topology costs and observable event counts. At minimum record
per-pass $(a,k,b,n)$, scan visits, activation frequency, executor transitions,
class-specific completions and service costs, and full pass/return intervals.
Do not fit a fresh unexplained residual for every lock. The present FC-PQ
\texttt{combiner\_time\_stat} reuses its start timestamp inside each delegate,
so it is not a full-pass timer; the benchmark's requester-equals-executor
flag is not a complete history of who performed combining work.

Two falsifiable regime predictions follow: at matched work and topology,
exposed per-operation migration decreases with $F/B=p/b$ until another
bottleneck dominates; sparse FC activity tracks retained-list visits per
completion, while FC-PQ tracks activation/pop work per completion rather than
only $N$. Neither prediction assumes a service-fairness bound.

Hold out footprints, operation mixes, and placements after calibration, then
predict time and crossover direction. Report uncertainty and failures,
including locality effects outside the repeated-footprint assumption.
Platform-specific ownership events and controlled access patterns are needed
to attribute coherence; generic LLC misses cannot identify protected-state
versus publication/PQ traffic. No coefficients or machine measurements are
fabricated here. The \texttt{crossover} and \texttt{batch\_frontier} artifact
checks exercise the algebra, not hardware speed or individual starvation.
